\part{Decisions}\label{tm:part:decisions}

\chapter{Decision log}\label{tm:ch:decision-log}

This chapter holds the design decisions. Each section gives the question, the options and the result.

\begin{information}
A decision record states the facts of its date. Some records state that Phase 5 was signed off at an ASAP7 frequency. This statement does not stand. Read \cref{tm:ch:standing-caveats} before you quote a figure from a decision record.
\end{information}

\begin{relocated}{docs/3PLL_CDC_EVALUATION.md}
% Converted from docs/3PLL_CDC_EVALUATION.md by tools/docs/md2tex.py; edit this file, not the original.
\section{3-PLL Multi-Clock + CDC Evaluation (pre-Phase-6 \#4)}\label{tm:dec:3pll-cdc}

\textbf{Status:} DECIDED 2026-06-27 --- \textbf{GO, 2-domain MVP}
(Option A): one CPU domain @ \textasciitilde1282 MHz + one shared
GPU+bus+peripheral domain @ 571 MHz, with a single async-AXI CDC bridge
on CPU↔crossbar. The full 3-domain split (separate GPU PLL) is rejected
as over-scoped (GPU already == bus speed). Folds in the \code{ste}
PLL-enable fix. Implementation = per-module PRs under the epic below.
\textbf{Date:} 2026-06-27 · Bead \code{claude_verilog_test-d8v}.
\textbf{Proposal:} instantiate the \code{pll_subsystem} 3× --- separate
PLLs for CPU / GPU / SoC-bus+peripherals --- to create genuine clock
domains, and evaluate whether that\textquotesingle s meaningful for
exercising CDC.

\subsection{Where we are today}\label{tm:dec:3pll-cdc:1-where-we-are-today}\origsec{1}

\begin{itemize}
\tightlist
\item
  \textbf{Single clock.} Every instance runs on \code{core_clk} from one
  \code{pll_subsystem} (\code{clk_i} ref); only \code{pll_subsystem}
  itself is on \code{clk_i} (bootstrap rule). The AXI4 crossbar /
  AXI-Lite ring / APB subtree are \textbf{purely synchronous} ---
  \code{soc_top.sv}, \code{soc_bus.sv}.
\item
  \textbf{Zero genuine CDC.} The only synchronizers are 2-FF on async
  \emph{pads} (IRQ \code{irq_src_i}, UART RX). No async-AXI FIFO, no
  gray-code pointers, no CDC-aware fabric.
  \texttt{design\_\allowbreak{}state.\allowbreak{}json:\allowbreak{}\ cdc\_\allowbreak{}clean=\allowbreak{}false};
  \code{cdc_snitch} reports 0 real internal crossings (BAD count is
  \textasciitilde95\% reset-domain false positives).
\item
  \textbf{Signed off single-clock} at M11: \textbf{571 MHz,
  GPU-governed} (\code{pnr/constraints/.../phase5_soc.sdc} = one
  \code{create_clock}, no
  \texttt{set\_\allowbreak{}clock\_\allowbreak{}groups\ -\allowbreak{}asynchronous}).
  The plan already flags multi-clock as a \textbf{High/Med risk} (CDC
  review; pre-harden macros {[}done{]}; per-domain SDC).
\end{itemize}

\subsection{The upside (why 3-PLL was proposed)}\label{tm:dec:3pll-cdc:2-the-upside-why-3-pll-was-proposed}\origsec{2}

{
\begin{longtable}[]{@{}lll@{}}
\toprule\noalign{}
Block & Standalone ASAP7 fmax & In single-clock SoC \\
\midrule\noalign{}
\endhead
\bottomrule\noalign{}
\endlastfoot
CPU & \textbf{1282 MHz} (780 ps, +24 ps; M7) & throttled to
\textbf{571} \\
GPU & \textbf{571 MHz} (its native fmax) & 571 (governs) \\
Bus/\allowbreak{}periph & fast (small logic) & 571 \\
\end{longtable}
}

So independent domains could run \textbf{CPU \textasciitilde1282 MHz
(2.25×)} while GPU stays 571. Plus: it\textquotesingle s the
\textbf{only way to exercise real CDC} in this design (currently none).

\subsection{The cost (this is the hard part)}\label{tm:dec:3pll-cdc:3-the-cost-this-is-the-hard-part}\origsec{3}

A domain split puts a \textbf{clock-domain crossing on every fabric
boundary where the master and the crossbar differ}. From the topology
map:

{
\begin{longtable}[]{@{}
  >{\raggedright\arraybackslash}p{\dimexpr(\linewidth-4\tabcolsep)*352/1000\relax}
  >{\raggedleft\arraybackslash}p{\dimexpr(\linewidth-4\tabcolsep)*113/1000\relax}
  >{\raggedright\arraybackslash}p{\dimexpr(\linewidth-4\tabcolsep)*515/1000\relax}@{}}
\toprule\noalign{}
\begin{minipage}[b]{\linewidth}\raggedright
Crossing
\end{minipage} & \begin{minipage}[b]{\linewidth}\raggedleft
Signals
\end{minipage} & \begin{minipage}[b]{\linewidth}\raggedright
Needs
\end{minipage} \\
\midrule\noalign{}
\endhead
\bottomrule\noalign{}
\endlastfoot
CPU ↔ crossbar (M0) & \textasciitilde70 & async-AXI bridge (FIFO,
gray-code ptrs) \\
GPU-data ↔ crossbar (M2) & \textasciitilde70 & async-AXI bridge \\
GPU-ifetch ↔ crossbar (M1) & \textasciitilde40 & async-AXI bridge
(read-only) \\
GPU-ctrl AXI-Lite & \textasciitilde20 & CDC on the AXI-Lite slave \\
(DMA, if its own domain) & \textasciitilde70 & async-AXI bridge \\
\end{longtable}
}

Plus: \textbf{reset synchronizers} per domain (reset is single-domain
today), \textbf{multi-clock SDC} (3× \code{create_clock} +
\texttt{set\_\allowbreak{}clock\_\allowbreak{}groups\ -\allowbreak{}asynchronous}
+ \code{max_delay} on every CDC path), \textbf{CDC verification} (the
\code{cdc_snitch} POC needs a wrapper --- reset-domain neutralization,
port allowlist, \code{magic_cdc} attributes --- before \code{BAD=0} is a
real gate; Verilator CDC is weak; metastability isn\textquotesingle t
fully simulatable), and a \textbf{full M11-scale PD re-closure} (3 CTS
trees, per-domain PDN/IR, multi-clock STA).

Async-AXI FIFOs / gray-code CDC are the \textbf{classic silicon-killer
RTL} --- hardest to get right, hardest to verify. This is a large,
high-risk feature, not a refactor.

\subsection{Key scoping insight --- the full 3-domain split is
over-scoped}\label{tm:dec:3pll-cdc:4-key-scoping-insight-the-full-3-domain-split-is-over-scoped}\origsec{4}

The proposal was CPU / GPU / bus as \textbf{three} domains. But
\textbf{GPU\textquotesingle s fmax (571) == the bus target (571)} --- so
GPU and bus can share one domain. That collapses the design to
\textbf{two} domains:

\begin{itemize}
\tightlist
\item
  \textbf{Domain A --- CPU @ \textasciitilde1282 MHz} (its own PLL).
\item
  \textbf{Domain B --- GPU + bus + peripherals @ 571 MHz} (one PLL).
\end{itemize}

Result: \textbf{ONE CDC boundary (CPU ↔ crossbar)} instead of four,
while still (a) exercising real CDC end-to-end and (b) delivering the
entire CPU 2.25× speedup. A separate 3rd PLL for GPU buys a
\emph{second} CDC boundary for \textasciitilde no benefit (GPU is
already at bus speed). \textbf{3 PLLs ≠ 3 useful domains here.}

\subsection{Recommendation}\label{tm:dec:3pll-cdc:5-recommendation}\origsec{5}

\textbf{3-PLL as literally proposed (3 domains): NO.} \textbf{A 2-domain
MVP (CPU-fast + everything-else): the right scope IF we pursue CDC} ---
and it\textquotesingle s worth doing for the CDC-methodology value (this
SoC has never exercised real CDC) plus a free CPU 2.25×, at a fraction
of the 3-domain cost/\allowbreak{}risk.

But note the honest counter: \textbf{Phase 5 is signed off at 571 MHz
single-clock, and this is a GPU-governed SoC} --- the CPU is the control
core, and CoreMark-class firmware (M10) runs fine. So multi-clock is
\textbf{not required to ship}; its value is (a) CDC
methodology/\allowbreak{}portfolio demonstration and (b) CPU-bound
headroom that this design doesn\textquotesingle t centrally need.

\subsubsection{Decision options (this PR)}\label{tm:dec:3pll-cdc:decision-options-this-pr}

\begin{itemize}
\tightlist
\item
  \textbf{A --- 2-domain MVP (recommended if pursuing CDC):} one
  async-AXI bridge on CPU↔crossbar; CPU PLL @1282, shared GPU+bus PLL
  @571; multi-clock SDC; wrap \code{cdc_snitch} to a real gate; folds in
  the \code{ste} PLL-enable fix (PLL control redesigned anyway). Own
  feature epic + re-verify + a PD re-closure.
\item
  \textbf{B --- Full 3-domain:} as proposed; adds the GPU↔fabric CDC
  boundary. More cost/risk for marginal benefit.
\item
  \textbf{C --- NO-GO (defer):} stay single-clock (signed off).
  Multi-clock/\allowbreak{}CDC isn\textquotesingle t needed for this
  GPU-centric SoC; revisit only if a CPU-bound workload justifies it.
  Keep \code{ste} as a standalone PLL-safety fix.
\end{itemize}

\subsubsection{Folds in}\label{tm:dec:3pll-cdc:folds-in}

The 3-PLL/\allowbreak{}2-domain work redesigns PLL control (multiple
\code{pll_subsystem}), so \textbf{fix the \code{pll_enable} self-brick
(\code{ste} / issue \#89) here} --- make CONTROL{[}0{]} safe per domain.

\subsection{If GO (A) --- implementation sketch (own PRs, each
re-verified)}\label{tm:dec:3pll-cdc:6-if-go-a-implementation-sketch-own-prs-each-re-verified}\origsec{6}

\begin{enumerate}
\def\labelenumi{\arabic{enumi}.}
\tightlist
\item
  \textbf{\code{async_axi_fifo.sv}} --- parameterized dual-clock AXI4
  CDC bridge (gray-code FIFO per channel). The crux; unit-test
  exhaustively (full-throttle, backpressure,
  near-empty/\allowbreak{}full, reset-while-busy).
\item
  \textbf{3×/2× \code{pll_subsystem}} in soc\_top (params ready) +
  \code{ste} PLL-enable safety fix.
\item
  Insert the async bridge on CPU↔crossbar; reset synchronizers; GPU+bus
  on the shared PLL (no CDC there).
\item
  \textbf{\code{phase5_soc_multiclock.sdc}} --- per-domain clocks,
  \texttt{set\_\allowbreak{}clock\_\allowbreak{}groups\ -\allowbreak{}asynchronous},
  \code{max_delay} CDC budgets.
\item
  \textbf{CDC verification}: wrap \code{cdc_snitch} (reset neutralize +
  port allowlist + \code{magic_cdc} on synchronizer FFs) → \code{BAD=0}
  CI gate; directed CPU↔fabric CDC tests; \code{soc_all} regression.
\item
  \textbf{PD re-closure}: multi-clock M11 run (2 CTS trees, per-domain
  STA/PDN).
\end{enumerate}

\subsection{Verdict line}\label{tm:dec:3pll-cdc:7-verdict-line}\origsec{7}

Pending human GO/NO-GO. Recommendation: \textbf{A (2-domain MVP)} if the
goal includes real-CDC methodology; \textbf{C (defer)} if the goal is
shipping the GPU-centric SoC, since single-clock 571 MHz is already
signed off.
\end{relocated}

\begin{relocated}{docs/M10_L2_DECISION_ANALYSIS.md}
% Converted from docs/M10_L2_DECISION_ANALYSIS.md by tools/docs/md2tex.py; edit this file, not the original.
\section{Phase 5 M10 --- L2 Cache Decision Analysis}\label{tm:dec:m10-l2}

\textbf{Status:} FINAL 2026-06-21 --- \textbf{NO-GO on both L2 and
2-way; keep the direct-mapped L1 D\$} (human architectural sign-off).
M11 floorplan is L2-free; the parameterized 2-way RTL stays on branch
\code{experiment/dcache-2way} as a record and is NOT merged. See §10 for
the final decision and its full basis. \textbf{Date:} 2026-06-21
\textbf{Gate criterion}
(\texttt{Doc:\allowbreak{}\ ds:\allowbreak{}ch:\allowbreak{}phase5-\allowbreak{}soc-\allowbreak{}integration-\allowbreak{}plan}
M10; \src{CLAUDE.md} locked-decision \#4):

\begin{quote}
Add \code{rtl/mem/l2_cache.sv} \textbf{only if} Phase 5 benchmarking
shows L1 miss rates are a bottleneck. Document the decision either way.
\end{quote}

\subsection{What was measured}\label{tm:dec:m10-l2:1-what-was-measured}\origsec{1}

M9 recorded cycle counts but not miss rates. M10 built a real RISC-V
bare-metal benchmark path (\src{sw/bench/}, riscv32-none-elf-gcc) and a
SoC cocotb harness (\src{tb/cocotb/soc/test_l2_bench.py}) that runs C on
the full \code{soc_top}, dumps the M7 hardware perf counters
(\code{mhpmcounter3}=I\$ miss, \code{mhpmcounter4}=D\$ miss,
\code{mcycle}, \code{minstret}) through a flushed scratch block, and
computes miss rates.

\begin{quote}
\textbf{Note 2026-10-10 (bead \bead{kiit}, \ghref{260}):} \code{mcycle} and
the other \code{rv32i_csr_file} performance counters stopped counting in
every cycle a CSR instruction was in EX (and in trap-entry / MRET
cycles). Cycle counts taken from \code{mcycle} before the fix therefore
under-count by one per CSR instruction in the measured window (the size
of the effect depends on how many CSR instructions the harness executes
inside its window; it was not re-measured). The results below are left
as recorded. See Doc:
cr:fix:fixes-index:14-performance-counters-freeze-while-a-csr-instruction-is-in-ex-bead-kiit-gh-260.
\end{quote}

The decision workload is \code{sweep.c}: a working-set sweep (512 B → 32
KB) in two access patterns --- \textbf{SEQ} (sequential) and
\textbf{STRIDE} (256 B stride, defeats spatial locality) --- each warmed
by a prime pass then measured at steady state. L1 D\$ = 4 KB,
direct-mapped, 16 B lines.

Raw data: \src{tb/cocotb/soc/l2_bench_results.md}.

\subsection{Result --- a hard capacity cliff at the 4 KB L1 boundary}\label{tm:dec:m10-l2:2-result-a-hard-capacity-cliff-at-the-4-kb-l1-boundary}\origsec{2}

{
\begin{longtable}[]{@{}rrr@{}}
\toprule\noalign{}
working set & SEQ D\$miss/\allowbreak{}1k-instr & STRIDE
D\$miss/\allowbreak{}1k-instr \\
\midrule\noalign{}
\endhead
\bottomrule\noalign{}
\endlastfoot
≤ 4 KB & \textbf{0.00} & \textbf{0.00} \\
8 KB & \textbf{49.96} & \textbf{189.35} \\
16 KB & 49.98 & 194.53 \\
32 KB & 49.99 & 197.23 \\
\end{longtable}
}

\begin{itemize}
\tightlist
\item
  \textbf{Below/at 4 KB:} zero D\$ misses --- the working set fits
  entirely in L1.
\item
  \textbf{First size above L1 (8 KB):} miss rate snaps to
  \textasciitilde50/\allowbreak{}1k-instr (SEQ) and
  \textasciitilde190/\allowbreak{}1k-instr (STRIDE) and then
  \textbf{plateaus} --- a textbook capacity wall, not a gradual
  roll-off.
\item
  \textbf{STRIDE ≈ 3.8× worse than SEQ} above the cliff --- classic
  direct-mapped conflict/\allowbreak{}eviction behavior.
\item
  \textbf{I\$ misses ≈ 0} for all sizes ≥ 1 KB --- the hot loops fit the
  4 KB I\(; the
  512 B/SEQ I\) figure is just cold first-touch. \textbf{I\$ is not a
  concern.}
\item
  Corroboration: a 48×48 int \code{matmul} (27 KB working set, 6.75× L1)
  showed \textasciitilde50 D\$miss/\allowbreak{}1k-instr before its
  software-multiply runtime made the sim intractable --- consistent with
  the SEQ plateau.
\end{itemize}

\subsection{Interpretation}\label{tm:dec:m10-l2:3-interpretation}\origsec{3}

\textbf{What the data proves:} the L1 D\$ has a sharp 4 KB capacity
boundary, and any data working set above it sustains a high miss rate to
SRAM (full AXI4 burst penalty per miss). A 16--32 KB L2 would absorb the
8 KB / 16 KB cases entirely and roughly halve the 32 KB SEQ miss stream.

\textbf{What the data does NOT prove:} that \emph{representative} SoC
CPU workloads actually operate above 4 KB with reuse. \code{sweep} is a
synthetic adversary built to defeat the cache --- it measures the
cache\textquotesingle s limit, not that real firmware reaches it. The
representative integer workloads (Dhrystone/\allowbreak{}CoreMark) were
not run to completion (no hardware multiply → software \code{__mulsi3}
makes them very slow in RTL sim).

\textbf{Architectural context} weighs against a reflexive "yes":

\begin{itemize}
\tightlist
\item
  This CPU is the \textbf{control/\allowbreak{}orchestration core} of a
  GPU-governed SoC (SoC fmax ≈ 571 MHz, GPU-bound). Bulk data-parallel
  work runs on the \textbf{GPU}, which has its own 16 KB shared memory +
  memory coalescer --- it does not lean on the CPU L1.
\item
  An L2 adds latency, area, and an \textbf{extra hard macro to the M11
  floorplan} (changing area/\allowbreak{}PDN/\allowbreak{}timing scope),
  against the conservative Phase 5 locked posture (decisions
  \#1/\#4/\#5: add capacity only when a real bottleneck is proven).
\item
  The STRIDE ≈ 4× penalty is a \textbf{conflict-miss} signature. If a
  cheaper lever is wanted first, a \textbf{2-way set-associative D\$}
  (CLAUDE.md decision \#2) attacks the conflict component directly and
  is far smaller than an L2.
\end{itemize}

\subsection{Recommendation}\label{tm:dec:m10-l2:4-recommendation}\origsec{4}

\textbf{Decision (2026-06-21): defer L2; evaluate a 2-way
set-associative D\$ first.}

L2 is not added in the Phase 5 base scope on this evidence alone.
Instead, attack the cheaper, more-targeted lever first: the STRIDE ≈ 4×
penalty is a \textbf{conflict-miss} signature (CLAUDE.md decision
\#2\textquotesingle s trigger --- "measurable conflict-miss impact" ---
is now met). A 2-way set-associative D\$ is far smaller than an L2 and
addresses the conflict component directly.

\textbf{Experiment plan (M10b):}

\begin{enumerate}
\def\labelenumi{\arabic{enumi}.}
\tightlist
\item
  Prototype a 2-way set-associative D\$ (\code{rv32i_dcache.sv}),
  parameterized ways, \textbf{same 4 KB total capacity} (128 sets × 2
  ways × 16 B) so the change isolates \emph{conflict}-miss reduction
  from capacity (capacity is held constant).
\item
  Re-run the same \code{sweep} benchmark and compare against the
  direct-mapped baseline in §2. Expected: SEQ plateau
  (\textasciitilde50/kI ≥8 KB) is \textbf{capacity} --- associativity
  will NOT move it; STRIDE (\textasciitilde190/kI) should drop toward
  SEQ if conflicts dominate it.
\item
  Re-decide L2 with that data: if 2-way collapses STRIDE and the
  residual SEQ capacity miss is deemed acceptable for control-plane
  firmware → no L2. If the SEQ capacity wall is judged a real bottleneck
  for on-core workloads → escalate to L2.
\end{enumerate}

This keeps the M11 floorplan L2-free for now; the 2-way D\$ (if adopted)
is a modification of the existing L1, not a new macro.

\begin{quote}
Direct-mapped remains the committed default until the 2-way experiment
justifies a change (Phase 3 locked decision \#2). The prototype runs on
an experiment branch.
\end{quote}

\subsection{2-Way D\$ Re-benchmark (M10b-2, 2026-06-21)}\label{tm:dec:m10-l2:5-2-way-d-re-benchmark-m10b-2-2026-06-21}\origsec{5}

\textbf{Correctness gate (Step 1):} RTL fix
\code{fr_null_20260621_125041_00} confirmed on branch
\code{experiment/dcache-2way} (commit \code{aba0c18} --- writeback
buffer now reads \code{cand_vw}, not stale \code{victim_way_q}).

\begin{itemize}
\tightlist
\item
  \texttt{make\ -\allowbreak{}C\ sim\ dcache}: TESTS=8 PASS=8 FAIL=0
  (including \code{test_dirty_eviction_writeback})
\item
  \texttt{make\ -\allowbreak{}C\ sim\ soc\_\allowbreak{}coherency}:
  TESTS=3 PASS=3 FAIL=0
\end{itemize}

\textbf{Benchmark (Step 2):} same \code{sweep} firmware, same
\code{test_l2_bench.py} harness, same \code{SRAM_MEM_WORDS=65536}.
Cache: \code{DCACHE_WAYS=2}, 128 sets × 2 ways × 16 B = 4 KB.

Raw data: \src{tb/cocotb/soc/l2_bench_results_2way.md}.

\subsubsection{Comparison table --- direct-mapped vs 2-way (same 4
KB capacity)}\label{tm:dec:m10-l2:51-comparison-table-direct-mapped-vs-2-way-same-4-kb-capacity}\origsec{5.1}

{
\begin{longtable}[]{@{}
  >{\raggedleft\arraybackslash}p{\dimexpr(\linewidth-10\tabcolsep)*115/1000\relax}
  >{\raggedleft\arraybackslash}p{\dimexpr(\linewidth-10\tabcolsep)*173/1000\relax}
  >{\raggedleft\arraybackslash}p{\dimexpr(\linewidth-10\tabcolsep)*173/1000\relax}
  >{\raggedleft\arraybackslash}p{\dimexpr(\linewidth-10\tabcolsep)*173/1000\relax}
  >{\raggedleft\arraybackslash}p{\dimexpr(\linewidth-10\tabcolsep)*173/1000\relax}
  >{\raggedright\arraybackslash}p{\dimexpr(\linewidth-10\tabcolsep)*173/1000\relax}@{}}
\toprule\noalign{}
\begin{minipage}[b]{\linewidth}\raggedleft
size (B)
\end{minipage} & \begin{minipage}[b]{\linewidth}\raggedleft
1-way SEQ /kI
\end{minipage} & \begin{minipage}[b]{\linewidth}\raggedleft
2-way SEQ /kI
\end{minipage} & \begin{minipage}[b]{\linewidth}\raggedleft
1-way STR /kI
\end{minipage} & \begin{minipage}[b]{\linewidth}\raggedleft
2-way STR /kI
\end{minipage} & \begin{minipage}[b]{\linewidth}\raggedright
STRIDE change
\end{minipage} \\
\midrule\noalign{}
\endhead
\bottomrule\noalign{}
\endlastfoot
512 & 0.00 & 0.00 & 0.00 & 0.00 & no change \\
1024 & 0.00 & 0.00 & 0.00 & 0.00 & no change \\
2048 & 0.00 & 0.00 & 0.00 & 0.00 & no change \\
4096 & 0.00 & 0.00 & 0.00 & 0.00 & no change \\
8192 & 49.96 & 49.96 & 189.35 & 189.35 & 0\% reduction \\
16384 & 49.98 & 49.98 & 194.53 & 194.53 & 0\% reduction \\
32768 & 49.99 & 49.99 & 197.23 & 197.23 & 0\% reduction \\
\end{longtable}
}

\subsubsection{Capacity vs conflict decomposition}\label{tm:dec:m10-l2:52-capacity-vs-conflict-decomposition}\origsec{5.2}

The two benchmarks are \textbf{bit-identical}. This outcome is
physically expected once the STRIDE access geometry is analyzed:

\begin{itemize}
\tightlist
\item
  The \code{sweep} firmware uses
  \texttt{STRIDE\_\allowbreak{}STEP\ =\allowbreak{}\ 64\ words\ =\allowbreak{}\ 256\ bytes}.
\item
  The D\$ line size is 16 bytes, and the cache has 256 sets (WAYS=1) or
  128 sets (WAYS=2).
\item
  A 256 B stride increments the cache-set index by 256/16 = 16 sets
  (WAYS=1) or 8 sets (WAYS=2). Each consecutive STRIDE access maps to a
  \emph{different set} in both configurations.
\item
  Because no two STRIDE accesses map to the same set, \textbf{there is
  no set aliasing}, and therefore \textbf{no conflict misses} --- only
  capacity misses (the working set exceeds 4 KB total cache space).
\item
  2-way associativity prevents evictions when two addresses map to the
  \emph{same set} with different tags. When no two live accesses share a
  set (as here), extra ways provide no benefit.
\end{itemize}

\textbf{Conclusion:} the STRIDE ≈ 3.8× penalty over SEQ measured in the
direct-mapped baseline is a \textbf{capacity miss phenomenon, not a
conflict-miss phenomenon}. The STRIDE pattern in this benchmark always
evicts by capacity; reducing the effective number of sets from 256 to
128 (the price of adding a second way at constant capacity) does not
change the miss rate because no aliasing was occurring in the first
place.

\subsubsection{Residual capacity wall}\label{tm:dec:m10-l2:53-residual-capacity-wall}\origsec{5.3}

Both WAYS=1 and WAYS=2 sustain \textasciitilde50 D\$miss/kI (SEQ) for
working sets ≥ 8 KB. This is a pure capacity wall: once the working set
exceeds 4 KB, every line brought in evicts a line that will be needed
again. The only remedy is more cache capacity (a larger L1 or an L2),
not more associativity.

The STRIDE ceiling (\textasciitilde190--197/\allowbreak{}kI) is
similarly capacity-driven. With a 256 B stride and a 16 B line, each
access is guaranteed cold regardless of cache size below the full
working-set footprint. Neither associativity nor an L2 would help STRIDE
unless the L2 were large enough to hold the entire working set.

\subsubsection{Implication for L2}\label{tm:dec:m10-l2:54-implication-for-l2}\origsec{5.4}

\begin{itemize}
\tightlist
\item
  \textbf{2-way D\$ at same 4 KB capacity: no measurable miss-rate
  benefit} for the \code{sweep} workload (which uses non-aliasing large
  strides).
\item
  \textbf{Residual capacity wall above 4 KB is entirely a capacity
  problem}, addressable only by a larger cache (a bigger L1 or an L2).
\item
  \textbf{An L2 of 16--32 KB would absorb the 8--16 KB working-set
  cases} (SEQ \textasciitilde50/kI would drop to \textasciitilde0;
  STRIDE would drop only if the L2 is large enough to hold the strided
  footprint --- typically not for STRIDE\_\allowbreak{}STEP
  \textless\textless{} working\_\allowbreak{}set).
\item
  Whether the capacity wall represents a \emph{real bottleneck} for the
  SoC\textquotesingle s actual control-plane firmware remains unproven
  (see §3 architectural context).
\end{itemize}

\begin{quote}
\textbf{L2 GO/NO-GO decision left to the human} on this data. The 2-way
experiment confirms: conflict misses are not the driver for this
workload; any remaining miss-rate concern is purely capacity, and only
an L2 (or larger L1) would address it.
\end{quote}

\subsection{Reproduce}\label{tm:dec:m10-l2:6-reproduce}\origsec{6}

\begin{lstlisting}[language=bash]
# Direct-mapped baseline (DCACHE_WAYS=1 in rv32i_cache_pkg.sv)
nix develop .#bench --command make -C sw/bench clean all   # build benchmarks
cd sim && make l2_bench                                      # writes l2_bench_results.md

# 2-way re-benchmark (branch experiment/dcache-2way, DCACHE_WAYS=2)
git checkout experiment/dcache-2way
cd sim && make l2_bench                                      # writes l2_bench_results_2way.md
\end{lstlisting}

\subsection{Conflict-stress benchmark (M10c, 2026-06-21)}\label{tm:dec:m10-l2:7-conflict-stress-benchmark-m10c-2026-06-21}\origsec{7}

\textbf{Purpose:} M10b\textquotesingle s \code{sweep} benchmark used a
256 B stride which never aliases the same cache set --- so it tested
\emph{capacity} misses only. To determine whether \textbf{conflict
misses} can be relevant for real SoC workloads, M10c (\code{conflict.c})
forces true set-aliasing: all K probe addresses are at 4096 B intervals
(same D\$ set for both WAYS=1 and WAYS=2). K ∈ \{1, 2, 4, 8, 16\}. K=2
is the key discriminator.

\subsubsection{Results --- D\$miss/kI by K-value}\label{tm:dec:m10-l2:71-results-dmisski-by-k-value}\origsec{7.1}

Branch \code{experiment/dcache-2way}. Same firmware binary
(\code{conflict.hex}) for both WAYS. ROUNDS=256 measured-pass accesses
per K-value.

{
\begin{longtable}[]{@{}
  >{\raggedleft\arraybackslash}p{\dimexpr(\linewidth-10\tabcolsep)*51/1000\relax}
  >{\raggedleft\arraybackslash}p{\dimexpr(\linewidth-10\tabcolsep)*152/1000\relax}
  >{\raggedleft\arraybackslash}p{\dimexpr(\linewidth-10\tabcolsep)*152/1000\relax}
  >{\raggedleft\arraybackslash}p{\dimexpr(\linewidth-10\tabcolsep)*169/1000\relax}
  >{\raggedleft\arraybackslash}p{\dimexpr(\linewidth-10\tabcolsep)*169/1000\relax}
  >{\raggedright\arraybackslash}p{\dimexpr(\linewidth-10\tabcolsep)*287/1000\relax}@{}}
\toprule\noalign{}
\begin{minipage}[b]{\linewidth}\raggedleft
K
\end{minipage} & \begin{minipage}[b]{\linewidth}\raggedleft
WAYS=1 D\$miss/kI
\end{minipage} & \begin{minipage}[b]{\linewidth}\raggedleft
WAYS=2 D\$miss/kI
\end{minipage} & \begin{minipage}[b]{\linewidth}\raggedleft
WAYS=1 miss/\allowbreak{}access
\end{minipage} & \begin{minipage}[b]{\linewidth}\raggedleft
WAYS=2 miss/\allowbreak{}access
\end{minipage} & \begin{minipage}[b]{\linewidth}\raggedright
verdict
\end{minipage} \\
\midrule\noalign{}
\endhead
\bottomrule\noalign{}
\endlastfoot
1 & \textbf{0.00} & \textbf{0.00} & 0.0000 & 0.0000 & fits both (no
conflict) \\
2 & \textbf{153.48} & \textbf{0.00} & \textbf{1.0000} & \textbf{0.0000}
& \textbf{KEY: 1-way thrashes, 2-way holds} \\
4 & 173.68 & 173.68 & 1.0000 & 1.0000 & both thrash (K \textgreater{} 2
ways) \\
8 & 185.91 & 185.91 & 1.0000 & 1.0000 & both thrash \\
16 & 192.70 & 192.70 & 1.0000 & 1.0000 & both thrash \\
\end{longtable}
}

Raw data: \src{tb/cocotb/soc/conflict_results_1way.md},
\src{tb/cocotb/soc/conflict_results_2way.md}.

\subsubsection{K=2 headline numbers}\label{tm:dec:m10-l2:72-k2-headline-numbers}\origsec{7.2}

{
\begin{longtable}[]{@{}lrrrr@{}}
\toprule\noalign{}
config & ΔD\$miss & D\$miss/\allowbreak{}access & D\$miss/kI & cycles
(measured pass) \\
\midrule\noalign{}
\endhead
\bottomrule\noalign{}
\endlastfoot
WAYS=1 & 512 & \textbf{1.0000} & 153.48 & 20 749 \\
WAYS=2 & 0 & \textbf{0.0000} & 0.00 & 16 141 \\
\end{longtable}
}

\textbf{K=2 WAYS=1:} every access evicts the other line → 1.0000
miss/\allowbreak{}access. The measured pass counts exactly 512 = 2 ×
ROUNDS misses --- one miss for each of the 512 accesses. \textbf{K=2
WAYS=2:} both lines reside simultaneously in the 2-way LRU set → 0
misses. Cycle counts: 20 749 (1-way, stall-heavy) vs 16 141 (2-way, no
stalls) --- a 22\% speedup for the K=2 conflict-heavy access pattern.

The 2-way LRU implementation (cand\_vw / lru\_array fix from
\code{fr_null_20260621_125041_00}) is \textbf{confirmed functionally
correct}.

\subsubsection{Interpretation}\label{tm:dec:m10-l2:73-interpretation}\origsec{7.3}

\textbf{What the conflict-stress benchmark proves:}

\begin{enumerate}
\def\labelenumi{\arabic{enumi}.}
\item
  \textbf{2-way D\$ eliminates conflict misses completely when K=2} ---
  the LRU policy works as designed; the cand\_vw RTL fix is confirmed.
\item
  \textbf{For K\textgreater2, both configs thrash equally} --- a 2-way
  cache cannot hold K\textgreater2 aliased lines any more than a 1-way
  cache can hold K\textgreater1. Adding associativity beyond 2 would
  help K=3 but not K≥4.
\item
  \textbf{The discriminating question} is whether real SoC firmware has
  hot data structures that map exactly 2 or more congruent lines to the
  same cache set. For this CPU (RV32I control core, not the
  data-parallel GPU):

  \begin{itemize}
  \tightlist
  \item
    Small scalar kernels (boot, peripheral drivers, interrupt handlers):
    working set typically \textless{} 512 B → all in L1, zero misses
    regardless of WAYS.
  \item
    \code{matmul}-class loops with two large arrays: both arrays may
    alias the same sets, making this a real concern. But the GPU handles
    bulk data-parallel work.
  \item
    Memory-scanning loops (boot ROM copy, DMA setup): stride-based, not
    aliasing.
  \end{itemize}
\item
  \textbf{The conflict-miss problem is real but workload-conditional.}
  The experiment proves 2-way works; it does not prove the workload
  reaches K=2 aliasing in practice.
\end{enumerate}

\subsubsection{Final architecture observation}\label{tm:dec:m10-l2:74-final-architecture-observation}\origsec{7.4}

{
\begin{longtable}[]{@{}
  >{\raggedright\arraybackslash}p{\dimexpr(\linewidth-2\tabcolsep)*484/1000\relax}
  >{\raggedright\arraybackslash}p{\dimexpr(\linewidth-2\tabcolsep)*496/1000\relax}@{}}
\toprule\noalign{}
\begin{minipage}[b]{\linewidth}\raggedright
axis
\end{minipage} & \begin{minipage}[b]{\linewidth}\raggedright
finding
\end{minipage} \\
\midrule\noalign{}
\endhead
\bottomrule\noalign{}
\endlastfoot
2-way eliminates conflict misses (K=2) & CONFIRMED --- 1.0000 → 0.0000
miss/\allowbreak{}access \\
Capacity wall (working set \textgreater{} 4 KB) & UNCHANGED --- requires
L2 or larger L1 \\
K≥4 conflict (\textgreater2 aliased hot lines) & UNSOLVED by 2-way;
needs ≥4-way or L2 \\
Real workload set aliasing & UNPROVEN --- sweep/\allowbreak{}conflict
are synthetic \\
\end{longtable}
}

\begin{quote}
\textbf{GO/NO-GO for 2-way promotion and/or L2 is left to the human.}
The conflict-stress result closes the open measurement question:
\textbf{2-way does eliminate true conflict misses when they exist}.
Whether the SoC CPU firmware encounters them is an
architectural-judgement call, not a simulation question.
\end{quote}

\subsection{Reproduce (M10c)}\label{tm:dec:m10-l2:8-reproduce-m10c}\origsec{8}

\begin{textblock}
# On branch experiment/dcache-2way (DCACHE_WAYS=2 active):
# conflict.hex is already built (sw/bench/build/conflict.hex, commit 87c16d1).
# Do NOT need to rebuild firmware — same binary for both WAYS.

# WAYS=1 run: edit rv32i_cache_pkg.sv to DCACHE_WAYS=1 then:
cd sim && make conflict_bench_1way     # → tb/cocotb/soc/conflict_results_1way.md

# WAYS=2 run: edit rv32i_cache_pkg.sv to DCACHE_WAYS=2 then:
cd sim && make conflict_bench_2way     # → tb/cocotb/soc/conflict_results_2way.md
\end{textblock}

\subsection{CoreMark Real-Workload Run (M10c-3)}\label{tm:dec:m10-l2:9-coremark-real-workload-run-m10c-3}\origsec{9}

\textbf{Branch}: \code{experiment/dcache-2way},
\textbf{DCACHE\_\allowbreak{}WAYS=1} (shipping direct-mapped, reverted
for this run). \textbf{Firmware}: \code{sw/bench/coremark} --- 3-kernel
port (list sort + 2×2 matmul + CRC-16), ITERATIONS=1.
\textbf{Working-set}: \textasciitilde232 B (list 80 B + matmul 24 B +
CRC buf 16 B + overhead) --- \textless{} 6\% of 4 KB L1 D\$.

{
\begin{longtable}[]{@{}lr@{}}
\toprule\noalign{}
metric & value \\
\midrule\noalign{}
\endhead
\bottomrule\noalign{}
\endlastfoot
Δinstret & 1,463 \\
Δcycles & 7,750 \\
ΔI\$-miss & 45 \\
\textbf{ΔD\$-miss} & \textbf{14} \\
I\$-miss / 1k-instr & 30.76 \\
\textbf{D\$-miss / 1k-instr} & \textbf{9.57} \\
D\$-miss / 1k-cycle & 1.81 \\
\end{longtable}
}

\textbf{Conclusion}: D\(-miss/kI = 9.57 — all 14 D\) misses are cold
fills on the first touch of the \textasciitilde232 B data region. After
warm-up the working set is fully resident in the 4 KB direct-mapped L1
with zero capacity or conflict pressure. Representative embedded
firmware does not drive L1 D\$ miss pressure; the L2 go/no-go decision
is driven by application working-set size (capacity cliff at 4-8 KB,
from the sweep benchmark), not by CoreMark-class workloads.

Results file: \src{tb/cocotb/soc/coremark_results.md}.

\subsection{Final Decision (2026-06-21) --- NO-GO on L2 and 2-way;
keep direct-mapped}\label{tm:dec:m10-l2:10-final-decision-2026-06-21-no-go-on-l2-and-2-way-keep-direct-mapped}\origsec{10}

Human architectural sign-off after the full three-leg investigation.

\subsubsection{Evidence summary}\label{tm:dec:m10-l2:evidence-summary}

{\small\setlength{\tabcolsep}{4pt}

{
\begin{longtable}[]{@{}
  >{\raggedright\arraybackslash}p{\dimexpr(\linewidth-6\tabcolsep)*97/1000\relax}
  >{\raggedright\arraybackslash}p{\dimexpr(\linewidth-6\tabcolsep)*78/1000\relax}
  >{\raggedright\arraybackslash}p{\dimexpr(\linewidth-6\tabcolsep)*402/1000\relax}
  >{\raggedright\arraybackslash}p{\dimexpr(\linewidth-6\tabcolsep)*402/1000\relax}@{}}
\toprule\noalign{}
\begin{minipage}[b]{\linewidth}\raggedright
Leg
\end{minipage} & \begin{minipage}[b]{\linewidth}\raggedright
Benchmark
\end{minipage} & \begin{minipage}[b]{\linewidth}\raggedright
Result
\end{minipage} & \begin{minipage}[b]{\linewidth}\raggedright
Reading
\end{minipage} \\
\midrule\noalign{}
\endhead
\bottomrule\noalign{}
\endlastfoot
Capacity & \code{sweep} & 0 D\$miss/kI ≤4 KB; \textasciitilde50 (SEQ) /
\textasciitilde190 (STRIDE) ≥8 KB, plateau & Hard 4 KB capacity wall ---
but a \textbf{synthetic} adversary \\
Conflict & \code{conflict} & K=2: direct-mapped 153/kI → \textbf{2-way
0/kI} (−22\% cyc); K≥4 both thrash & 2-way \textbf{does} eliminate true
conflict misses --- proven + RTL-correct --- but only ≤2 ways \\
Real workload & \code{coremark} & 232 B working set; \textbf{9.57
D\$miss/\allowbreak{}kI, all 14 cold-fill, 0 steady-state evictions} &
Representative firmware \textbf{fits L1} with 18× margin --- no capacity
or conflict pressure \\
\end{longtable}
}

}

\subsubsection{Decision}\label{tm:dec:m10-l2:decision}

\textbf{No L2, and do not adopt the 2-way D\$. The shipping L1 stays 4
KB direct-mapped (\code{DCACHE_WAYS=1}).}

Rationale:

\begin{itemize}
\tightlist
\item
  The L1 capacity wall (sweep) and the 2-way conflict benefit (conflict)
  only appear under \textbf{synthetic} access patterns. The
  representative workload (CoreMark) sits 18× inside the 4 KB L1 and
  shows only cold-fill misses --- real control-plane firmware does not
  reach either regime.
\item
  This CPU is the control/\allowbreak{}orchestration core of a
  \textbf{GPU-governed} SoC (\textasciitilde571 MHz). Bulk data-parallel
  work runs on the GPU (own 16 KB shared mem + coalescer); it does not
  lean on the CPU L1.
\item
  An L2 adds latency, area, and an extra \textbf{M11 floorplan macro}
  (area/\allowbreak{}PDN/\allowbreak{}timing) for no demonstrated
  real-workload benefit.
\item
  A 2-way D\$ adds a \textasciitilde100--200 ps hit-path way-select mux
  for conflict insurance that representative firmware never cashes in.
\item
  Consistent with Phase 5 locked decisions
  \#1/\allowbreak{}\#2/\allowbreak{}\#4/\allowbreak{}\#5 (add
  capacity/\allowbreak{}associativity only when a real bottleneck is
  proven). The trigger was tested and \textbf{not} met by representative
  workloads.
\end{itemize}

\subsubsection{Disposition}\label{tm:dec:m10-l2:disposition}

\begin{itemize}
\tightlist
\item
  \code{experiment/dcache-2way} is \textbf{retained as the record}: it
  holds the verified parameterized 2-way RTL (WAYS=1 bit-identical
  baseline; WAYS=2 functionally correct, dcache 8/8 + coherency 3/3) and
  the conflict/\allowbreak{}CoreMark benchmarks. It is \textbf{not
  merged}.
\item
  \textbf{Revisit trigger:} if a future on-core workload is shown to
  sustain \textgreater4 KB working sets (→ L2 / larger L1) or K≥2
  congruent-line hot sets (→ 2-way), re-open this analysis. The
  benchmark harness (\src{sw/bench/} + \code{test_l2_bench.py}) and the
  2-way RTL are ready to re-measure.
\item
  \textbf{M11 proceeds} with the 4 KB direct-mapped L1, L2-free.
\end{itemize}
\end{relocated}

\begin{relocated}{docs/PERIPHERAL_BUS_EVALUATION.md}
% Converted from docs/PERIPHERAL_BUS_EVALUATION.md by tools/docs/md2tex.py; edit this file, not the original.
\section{Peripheral Bus Evaluation --- AXI4-Lite vs APB (pre-Phase-6)}\label{tm:dec:peripheral-bus}

\textbf{Status:} RECOMMENDATION --- pending human sign-off (this PR is
the decision + plan; migration module-PRs link back here).
\textbf{Date:} 2026-06-24 · Bead \code{claude_verilog_test-ahe}.
\textbf{Why now:} Phase 6 adds \textbf{6 more register peripherals}
(GPIO/\allowbreak{}I2C/\allowbreak{}PWM/\allowbreak{}WDT/\allowbreak{}TRNG/\allowbreak{}AES-SHA)
on this bus. Settle the peripheral-bus standard before scaling, or re-do
them.

\subsection{Question}\label{tm:dec:peripheral-bus:1-question}\origsec{1}

The SoC control path is an AXI4-Lite ring
(\code{axi_lite_interconnect.sv}) fed from the AXI4 crossbar via
\code{axi4_to_axilite.sv}. Do the downstream peripherals actually need
AXI4-Lite, or is a simpler APB bus sufficient?

\subsection{Findings (evidence-based)}\label{tm:dec:peripheral-bus:2-findings-evidence-based}\origsec{2}

\subsubsection{The ring carries two different classes of slave}\label{tm:dec:peripheral-bus:the-ring-carries-two-different-classes-of-slave}

\code{soc_periph_map_pkg.sv} --- 7 slaves @
0x2000\_\allowbreak{}1000--7FFF (4 KB each):

{
\begin{longtable}[]{@{}
  >{\raggedright\arraybackslash}p{\dimexpr(\linewidth-4\tabcolsep)*378/1000\relax}
  >{\raggedright\arraybackslash}p{\dimexpr(\linewidth-4\tabcolsep)*261/1000\relax}
  >{\raggedright\arraybackslash}p{\dimexpr(\linewidth-4\tabcolsep)*342/1000\relax}@{}}
\toprule\noalign{}
\begin{minipage}[b]{\linewidth}\raggedright
Slave
\end{minipage} & \begin{minipage}[b]{\linewidth}\raggedright
Class
\end{minipage} & \begin{minipage}[b]{\linewidth}\raggedright
Bus need
\end{minipage} \\
\midrule\noalign{}
\endhead
\bottomrule\noalign{}
\endlastfoot
\textbf{GPU ctrl} (0x\ldots1000) & accelerator
command/\allowbreak{}config & AXI4-Lite-native (anchors the ring) \\
\textbf{DMA ctrl} (0x\ldots5000) & master config & AXI4-Lite
(it\textquotesingle s also an AXI4 master) \\
\textbf{PLL regs} (0x\ldots7000) & config & simple register file \\
\textbf{UART / SPI / Timer / IRQ}
(0x\ldots2/\allowbreak{}3/\allowbreak{}4/\allowbreak{}6000) &
\textbf{simple register peripherals} & \textbf{APB-class} \\
\end{longtable}
}

\subsubsection{The 4 simple peripherals are APB-equivalent}\label{tm:dec:peripheral-bus:the-4-simple-peripherals-are-apb-equivalent}

All of \code{uart_controller}, \code{spi_controller}, \code{timer},
\code{interrupt_controller} instantiate \code{axi_lite_register_bank}
and are plain memory-mapped register files. No outstanding-transaction
depth \textgreater1, no burst, no pipelined/\allowbreak{}concurrent R+W
exploited, no locked/\allowbreak{}atomic, responses always OKAY.
(\code{axi_lite_register_bank.sv:21} = depth-1 write + depth-1 read;
\code{axi_lite_interconnect.sv:13} = depth-1 per channel; AXI-Lite has
no bursts.)

\subsubsection{Which AXI4-Lite features are genuinely used}\label{tm:dec:peripheral-bus:which-axi4-lite-features-are-genuinely-used}

{\small\setlength{\tabcolsep}{4pt}

{
\begin{longtable}[]{@{}
  >{\raggedright\arraybackslash}p{\dimexpr(\linewidth-4\tabcolsep)*306/1000\relax}
  >{\raggedright\arraybackslash}p{\dimexpr(\linewidth-4\tabcolsep)*131/1000\relax}
  >{\raggedright\arraybackslash}p{\dimexpr(\linewidth-4\tabcolsep)*542/1000\relax}@{}}
\toprule\noalign{}
\begin{minipage}[b]{\linewidth}\raggedright
Feature
\end{minipage} & \begin{minipage}[b]{\linewidth}\raggedright
Used?
\end{minipage} & \begin{minipage}[b]{\linewidth}\raggedright
Note / APB cover
\end{minipage} \\
\midrule\noalign{}
\endhead
\bottomrule\noalign{}
\endlastfoot
Multiple outstanding (\textgreater1) & \textbf{No} & depth-1 everywhere
→ APB single-transfer equivalent \\
Burst & \textbf{No} & the bridge serializes AXI4 bursts to single beats
already \\
\textbf{Byte strobes (\code{wstrb})} & \textbf{Yes} & partial-word
writes (UART/SPI byte lanes) → \textbf{APB4 \code{pstrb}} covers this \\
Concurrent R/W channels & not exploited & single CPU master serializes;
only matters for a future 2nd control master \\
DECERR (unmapped) & partial & generated by the interconnect, not the
slaves → APB \code{pslverr} covers \\
ID, AxPROT/\allowbreak{}LOCK/\allowbreak{}CACHE/\allowbreak{}QOS, EXOKAY
& \textbf{No} & tied off / bookkeeping only \\
\end{longtable}
}

}

\textbf{The only real AXI-Lite feature the simple peripherals use is
byte-strobe partial writes --- covered by APB4 \code{pstrb}.}

\subsection{Recommendation --- hybrid: AXI-Lite fabric + APB
peripheral subtree}\label{tm:dec:peripheral-bus:3-recommendation-hybrid-axi-lite-fabric-apb-peripheral-subtree}\origsec{3}

\textbf{Adopt APB for the simple-register-peripheral subtree; keep
AXI4-Lite for the accelerator/\allowbreak{}master ports.}

\begin{itemize}
\tightlist
\item
  Add one \textbf{\code{axil_to_apb} bridge} off the AXI-Lite ring
  feeding an \textbf{APB4 peripheral bus} (APB4 for \code{pstrb}).
\item
  Move \textbf{UART / SPI / Timer / IRQ} --- and all Phase 6 peripherals
  (GPIO/\allowbreak{}I2C/\allowbreak{}PWM/\allowbreak{}WDT/\allowbreak{}TRNG/\allowbreak{}AES-SHA)
  --- onto that APB subtree behind an \code{apb_register_bank} (same
  register logic, APB ports).
\item
  \textbf{Keep AXI4-Lite} for \textbf{GPU ctrl} and \textbf{DMA ctrl}
  (they are accelerators/\allowbreak{}masters, AXI-Lite-native --- the
  reason the ring exists) and optionally \textbf{PLL}.
\end{itemize}

This is the textbook SoC topology (AXI/\allowbreak{}AXI-Lite fabric +
masters, single AXI-Lite→APB bridge to a slow-peripheral APB tree). It
is the right structure \textbf{before} Phase 6 drops in 6 APB-class
peripherals, and it keeps AXI-Lite exactly where it is justified.

\textbf{Net:} APB is sufficient (indeed the better fit) for the simple
peripherals; AXI4-Lite stays for GPU/DMA. Gains: simpler per-peripheral
slave (no AW/\allowbreak{}W-decouple FSM), 1-cycle-faster reads (no
bridge latch on that subtree), and Phase 6 peripherals become trivial
APB drop-ins. Cost: one bridge + per-module bus-port swap + re-verify.

\subsubsection{Alternative considered --- keep all AXI4-Lite}\label{tm:dec:peripheral-bus:alternative-considered-keep-all-axi4-lite}

Defensible (cohesion, zero churn, peripherals already verified
\code{soc_all} 82/82). Rejected as the primary recommendation because it
makes every Phase 6 peripheral carry an AXI-Lite register-bank FSM for
no functional benefit, and APB is the conventional fit for these slow
register IPs.

\subsection{Migration plan (if approved) --- per-module PRs linked to
this one}\label{tm:dec:peripheral-bus:4-migration-plan-if-approved-per-module-prs-linked-to-this-one}\origsec{4}

This PR is the \textbf{main/\allowbreak{}tracking PR}. Each step below
is its \textbf{own PR}, linked back here (bead
\code{claude_verilog_test-mdu} epic + sub-beads). Re-run \code{soc_all}
(verification-orchestrator) after each.

\begin{enumerate}
\def\labelenumi{\arabic{enumi}.}
\tightlist
\item
  \textbf{\code{apb4_register_bank.sv} + \code{axil_to_apb.sv} bridge}
  (new RTL + unit tests). PR-1.
\item
  \textbf{UART} → APB register bank (swap bus ports; keep
  FIFO/\allowbreak{}snoop logic). PR-2.
\item
  \textbf{SPI} → APB. PR-3.
\item
  \textbf{Timer} → APB. PR-4.
\item
  \textbf{Interrupt controller} → APB. PR-5.
\item
  \textbf{\code{soc_top} + \code{soc_periph_map_pkg}} --- insert
  \code{axil_to_apb} on the ring, move the 4 (+PLL?) slaves onto the APB
  subtree, rewire. PR-6.
\item
  \textbf{Full \code{soc_all} regression} + (optional) a 1M-cycle SoC
  stress re-run. PR-7.
\end{enumerate}

Each PR keeps AXI-Lite for GPU/DMA untouched. APB4 (\code{pstrb})
preserves the byte-strobe writes. Delegation: RTL →
rtl-design-orchestrator; sims → verification-orchestrator.

\subsection{Decision}\label{tm:dec:peripheral-bus:5-decision}\origsec{5}

GO / NO-GO on the hybrid migration is a human architectural call (this
PR). If GO → the \code{mdu} migration epic\textquotesingle s per-module
PRs proceed; if NO-GO → close \code{ahe}/\code{mdu} with "keep AXI-Lite
ring as-is," and Phase 6 peripherals stay AXI-Lite-native.
\end{relocated}

\begin{relocated}{docs/XEZIM_EVALUATION.md}
% Converted from docs/XEZIM_EVALUATION.md by tools/docs/md2tex.py; edit this file, not the original.
\section{\texorpdfstring{Xezim evaluation: secondary compile +
simulation tool (bead
\code{qd9l})}{Xezim evaluation: secondary compile + simulation tool (bead )}}\label{tm:dec:xezim}

\textbf{Date:} 2026-11-01. \textbf{Status:} evaluation only. No CI, RTL
or flow change; adoption is an open human decision (section 9).

\textbf{Verdict: do not add xezim to cocotb regression or CI now.} It
silently miscompiles an idiom used in this RTL (section 4), and stock
cocotb 1.9.2 cannot run on it (section 3). It is usable locally as a
cross-check for functional gate-level sim and X-probing, never as sole
evidence. \textbf{Separately, slang closes the use-before-declare CI gap
today with no new dependency} (section 5).

Tags: \textbf{measured} (run here, this session), \textbf{inferred},
\textbf{claimed} (xezim\textquotesingle s docs, not reproduced),
\textbf{reported} (user or repo docs, not run here). Measurements are
single runs unless stated. Raw logs live under
\code{/nobackup/xezim_eval/} (not in the repo, and may be wiped).

\subsection{Method}\label{tm:dec:xezim:1-method}\origsec{1}

Tool under test: xezim 0.11.4, git \code{9dcffd69-dirty} (local one-line
\code{Cargo.toml} edit), plain release build, no JIT. Baselines:
Verilator 5.048, sv2v 0.0.13.1, pyslang 11.0.0, cocotb 1.9.2 (repo pin)
and 2.0.1 (isolated install). RTL: \code{origin/main} at \code{87b7cf0}.

Each legality check needed a negative control, and each pass needed a
control able to fail. \textbf{VCS and Xcelium are not installed on this
host}: their figures are reported from
\texttt{Doc:\allowbreak{}\ um:\allowbreak{}flow:\allowbreak{}proprietary-\allowbreak{}sim}
(2026-10-24), plus the user\textquotesingle s confirmation from another
machine that they reject use-before-declare.

Phase 3, the full regression bake-off, \textbf{was not run}: the
plan\textquotesingle s stop gate tripped (section 3).

\subsection{Elaboration (measured)}\label{tm:dec:xezim:2-elaboration-measured}\origsec{2}

\code{soc_top} and \code{rv32i_cpu_top} elaborate with exit 0 on four
define arms (default, \code{SRAM_SKY130}, \code{SRAM_ASAP7},
\code{FPGA}). \code{--verbose} confirms real elaboration: 92 modules on
default \code{soc_top}.

{
\begin{longtable}[]{@{}
  >{\raggedright\arraybackslash}p{\dimexpr(\linewidth-6\tabcolsep)*327/1000\relax}
  >{\raggedright\arraybackslash}p{\dimexpr(\linewidth-6\tabcolsep)*163/1000\relax}
  >{\raggedright\arraybackslash}p{\dimexpr(\linewidth-6\tabcolsep)*151/1000\relax}
  >{\raggedright\arraybackslash}p{\dimexpr(\linewidth-6\tabcolsep)*339/1000\relax}@{}}
\toprule\noalign{}
\begin{minipage}[b]{\linewidth}\raggedright
Tool, \code{soc_top}
\end{minipage} & \begin{minipage}[b]{\linewidth}\raggedright
Wall
\end{minipage} & \begin{minipage}[b]{\linewidth}\raggedright
Peak RSS
\end{minipage} & \begin{minipage}[b]{\linewidth}\raggedright
Note
\end{minipage} \\
\midrule\noalign{}
\endhead
\bottomrule\noalign{}
\endlastfoot
xezim \code{--compile} & 0.46-0.81 s & 135-186 MB & \\
Verilator \code{--lint-only} & 2.4-4.1 s & 91-151 MB & \\
sv2v & \textasciitilde4.7 s & \textasciitilde655 MB & \\
pyslang full compilation & 0.2-0.34 s & 57-73 MB & Default, FPGA, macro
arms \\
\end{longtable}
}

{
\begin{longtable}[]{@{}
  >{\raggedright\arraybackslash}p{\dimexpr(\linewidth-8\tabcolsep)*518/1000\relax}
  >{\raggedright\arraybackslash}p{\dimexpr(\linewidth-8\tabcolsep)*88/1000\relax}
  >{\raggedright\arraybackslash}p{\dimexpr(\linewidth-8\tabcolsep)*198/1000\relax}
  >{\raggedright\arraybackslash}p{\dimexpr(\linewidth-8\tabcolsep)*88/1000\relax}
  >{\raggedright\arraybackslash}p{\dimexpr(\linewidth-8\tabcolsep)*88/1000\relax}@{}}
\toprule\noalign{}
\begin{minipage}[b]{\linewidth}\raggedright
Negative control
\end{minipage} & \begin{minipage}[b]{\linewidth}\raggedright
xezim
\end{minipage} & \begin{minipage}[b]{\linewidth}\raggedright
Verilator
\end{minipage} & \begin{minipage}[b]{\linewidth}\raggedright
sv2v
\end{minipage} & \begin{minipage}[b]{\linewidth}\raggedright
slang
\end{minipage} \\
\midrule\noalign{}
\endhead
\bottomrule\noalign{}
\endlastfoot
Missing module & reject & reject & \textbf{accept} & reject \\
Use before declare (\code{assign}, \code{always_comb}) & \textbf{accept}
& \textbf{accept} & \textbf{accept} & reject \\
Never-declared identifier & reject & reject & \textbf{accept} &
reject \\
Illegal syntax & reject & reject & reject & reject \\
Override of nonexistent parameter & reject & reject & reject & reject \\
Override of localparam & reject & reject & reject & reject \\
Too many positional overrides & \textbf{accept} & reject & reject &
reject \\
Same parameter overridden twice & \textbf{accept} & reject & reject &
reject \\
String into \texttt{int\ unsigned} parameter & accept & accept + warning
& accept & accept \\
Positive control, legal string override & accept & accept & accept &
accept \\
\end{longtable}
}

The string-into-integer row is probably legal SV (inferred), so
accepting it isn\textquotesingle t scored as a defect. Struct-select
reproducers D1-D3
(\texttt{Doc:\allowbreak{}\ cr:\allowbreak{}defect:\allowbreak{}synlig-\allowbreak{}struct-\allowbreak{}select-\allowbreak{}ainf})
give correct values on xezim and Verilator; this is no evidence about
Synlig.

\subsubsection{Can xezim\textquotesingle s misses be fixed? (source
read, nothing built)}\label{tm:dec:xezim:can-xezims-misses-be-fixed-source-read-nothing-built}

{\footnotesize\setlength{\tabcolsep}{3pt}

{
\begin{longtable}[]{@{}
  >{\raggedright\arraybackslash}p{\dimexpr(\linewidth-4\tabcolsep)*189/1000\relax}
  >{\raggedright\arraybackslash}p{\dimexpr(\linewidth-4\tabcolsep)*341/1000\relax}
  >{\raggedright\arraybackslash}p{\dimexpr(\linewidth-4\tabcolsep)*450/1000\relax}@{}}
\toprule\noalign{}
\begin{minipage}[b]{\linewidth}\raggedright
Miss
\end{minipage} & \begin{minipage}[b]{\linewidth}\raggedright
Location
\end{minipage} & \begin{minipage}[b]{\linewidth}\raggedright
Assessment (effort inferred)
\end{minipage} \\
\midrule\noalign{}
\endhead
\bottomrule\noalign{}
\endlastfoot
Duplicate named override & \code{xezim-parser/src/strict_check.rs}
\textasciitilde268 & Small: seen-set, hours \\
Excess positional & Same pass, or \code{src/elaborate.rs}
\textasciitilde35346 & Small: count check, hours; must count body
parameters right \\
Use before declare, read & \code{src/elaborate.rs:7137} &
\textbf{Deliberate omission:} the code flags writes only, as its
reference simulator is said to tolerate reads "irregularly". Opt-in flag
recommended \\
\end{longtable}
}

}

All belong in the separate \code{xezim-core} repo, so upstreaming is the
only sensible long-term route.

\subsection{cocotb (measured)}\label{tm:dec:xezim:3-cocotb-measured}\origsec{3}

{\footnotesize\setlength{\tabcolsep}{3pt}

{
\begin{longtable}[]{@{}
  >{\raggedright\arraybackslash}p{\dimexpr(\linewidth-6\tabcolsep)*112/1000\relax}
  >{\raggedright\arraybackslash}p{\dimexpr(\linewidth-6\tabcolsep)*286/1000\relax}
  >{\raggedright\arraybackslash}p{\dimexpr(\linewidth-6\tabcolsep)*291/1000\relax}
  >{\raggedright\arraybackslash}p{\dimexpr(\linewidth-6\tabcolsep)*291/1000\relax}@{}}
\toprule\noalign{}
\begin{minipage}[b]{\linewidth}\raggedright
Setup
\end{minipage} & \begin{minipage}[b]{\linewidth}\raggedright
Trivial DUT
\end{minipage} & \begin{minipage}[b]{\linewidth}\raggedright
\code{test_gpio}
\end{minipage} & \begin{minipage}[b]{\linewidth}\raggedright
Other
\end{minipage} \\
\midrule\noalign{}
\endhead
\bottomrule\noalign{}
\endlastfoot
Stock cocotb 1.9.2 & \textbf{FATAL}
\texttt{We\ are\ calling\ up\ again}, no \code{results.xml} &
\textbf{FATAL} & \\
1.9.2 + throwaway shim & 10 pass, both fail controls report FAIL &
18/18, sim time per test identical to Verilator; RTL mutant caught &
\code{test_command_queue} 3/3, pyuvm smoke 4/4,
\code{test_mbist_soc_slice} \textbf{2/3} \\
cocotb 2.0.1, no shim & 3 pass, fail control reports FAIL & Not run &
\code{test_functional_path_and_no_fault_run} same failure as shim \\
\end{longtable}
}

}

\begin{itemize}
\tightlist
\item
  Trigger: \code{cbValueChange} delivered while Python is inside
  \code{vpi_put_value}, for example a Python-driven clock. Timer-only
  tests and HDL-generated clocks pass.
\item
  \textbf{Which side violates the VPI contract is not established.}
  cocotb 2 tolerates it, so it isn\textquotesingle t a blanket xezim
  failure.
\item
  The shim is a byte-patched copy of cocotb\textquotesingle s VPI
  library plus \code{LD_PRELOAD}. Not acceptable for production.
\item
  Project cannot simply move to cocotb 2: \code{requirements.txt:9}
  records that 2.x removed the \code{units=} keyword used throughout
  \src{tb/}. Migration size not measured.
\item
  Throughput: \code{test_gpio} on xezim \textasciitilde1.5 s end to end,
  no build. Verilator \textasciitilde25 s build, test body 1.16 s vs
  xezim 1.25-1.37 s. \textbf{Unit-suite only; long-run throughput
  unmeasured.}
\end{itemize}

\subsection{Silent miscompile (measured, reproduced twice
independently)}\label{tm:dec:xezim:4-silent-miscompile-measured-reproduced-twice-independently}\origsec{4}

\begin{lstlisting}[language=Verilog]
function automatic logic f(input logic [7:0] a); return a == 8'h40; endfunction
logic [7:0] x = 8'h40;  logic one = 1, zero = 0;
wire a = ~f(x) || zero;   // correct 0, xezim 0.11.4 gives 1
wire b = ~f(x) && one;    // correct 0, xezim 0.11.4 gives 1
\end{lstlisting}

{\small\setlength{\tabcolsep}{4pt}

{
\begin{longtable}[]{@{}
  >{\raggedright\arraybackslash}p{\dimexpr(\linewidth-2\tabcolsep)*614/1000\relax}
  >{\raggedright\arraybackslash}p{\dimexpr(\linewidth-2\tabcolsep)*366/1000\relax}@{}}
\toprule\noalign{}
\begin{minipage}[b]{\linewidth}\raggedright
Form
\end{minipage} & \begin{minipage}[b]{\linewidth}\raggedright
xezim
\end{minipage} \\
\midrule\noalign{}
\endhead
\bottomrule\noalign{}
\endlastfoot
\texttt{\textasciitilde{}f(x)\ ||\ y},
\texttt{\textasciitilde{}f(x)\ \&\&\ y} (wire, \code{always_comb}, NBA)
& \textbf{wrong} (\texttt{\&\&} measured on wire only) \\
\texttt{\textasciitilde{}f(x)} alone,
\texttt{\textasciitilde{}fx\ ||\ y} (no call), \texttt{!f(x)\ ||\ y},
\texttt{\textasciitilde{}f(x)\ |\ y} & right \\
\texttt{if\ (\textasciitilde{}f(x))},
\texttt{\textasciitilde{}f(x)\ ?\ a\ :\allowbreak{}\ b},
\texttt{1'(\textasciitilde{}f(x))\ ||\ y} & right \\
Same in \code{initial}, blocking & right \\
\end{longtable}
}

}

\begin{itemize}
\tightlist
\item
  Verilator gives correct results. \code{XEZIM_TWO_STATE=0},
  \code{XEZIM_JIT=0}, \code{XEZIM_INIT_ZERO=1} don\textquotesingle t
  change it.
\item
  Inferred mechanism, source not traced: \texttt{\textasciitilde{}}
  widened in logical-operand context, making the result always nonzero.
\item
  \textbf{This RTL hits it:}
  \texttt{rtl/\allowbreak{}soc/\allowbreak{}sram\_\allowbreak{}controller.\allowbreak{}sv:\allowbreak{}253-\allowbreak{}254,\allowbreak{}\ 647-\allowbreak{}648,\allowbreak{}\ 825-\allowbreak{}826},
  the only
  \texttt{\textasciitilde{}call(.\allowbreak{}.\allowbreak{}.\allowbreak{})}
  sites in \src{rtl/}. On xezim a legal SRAM write returns SLVERR.
  \textbf{Not an RTL bug} on this evidence: Verilator correct, and the
  RTL is legal.
\item
  It shows why xezim can\textquotesingle t be sole evidence: a similar
  miscompile that hides a fault instead of raising one would pass
  silently. Other forms were not swept.
\end{itemize}

\subsection{Slang as use-before-declare check (measured)}\label{tm:dec:xezim:5-slang-as-use-before-declare-check-measured}\origsec{5}

Gap: Verilator and sv2v accept use-before-declare, while VCS and Xcelium
reject it (reported), and neither runs in CI. \code{pyslang==11.0.0} is
\textbf{already pinned} in \src{requirements.txt} and used by the
blocking \code{check_upf.py} gate.

{\small\setlength{\tabcolsep}{4pt}

{
\begin{longtable}[]{@{}
  >{\raggedright\arraybackslash}p{\dimexpr(\linewidth-2\tabcolsep)*490/1000\relax}
  >{\raggedright\arraybackslash}p{\dimexpr(\linewidth-2\tabcolsep)*490/1000\relax}@{}}
\toprule\noalign{}
\begin{minipage}[b]{\linewidth}\raggedright
Check
\end{minipage} & \begin{minipage}[b]{\linewidth}\raggedright
Result
\end{minipage} \\
\midrule\noalign{}
\endhead
\bottomrule\noalign{}
\endlastfoot
Negative controls & Every illegal case rejected; legal cases accepted
(section 2) \\
\code{soc_top} default (396 instances, 84 unique modules), \code{-DFPGA}
& 0 errors, 0 warnings \\
\code{rv32i_cpu_top}, \code{gpu_top} default & 0 errors, 0 warnings \\
\code{soc_top} Sky130 + ASAP7 \code{__pnr__} arms & 0 errors with
\code{--ignore-unknown-modules} \\
Seeded use-before-declare: \code{soc_top}, \code{gpio_controller},
\code{rv32i_regfile} & All 3 caught, default arm; \code{soc_top} also on
Sky130 arm \\
\end{longtable}
}

}

Limits:

\begin{itemize}
\tightlist
\item
  \textbf{\code{--ignore-unknown-modules} hides missing modules}
  (measured: missing-module control exits 0). Macro arms need it because
  macro models aren\textquotesingle t in the file lists. Use it only
  there, or add macro stubs.
\item
  Testbench code wasn\textquotesingle t checked. Default warning set
  only; stricter \code{-Wextra} not evaluated.
\item
  This section reports the trial. The gate that followed is
  \src{tools/verif/check_elab_strict.py}, clean on its 6 arms: 3
  defaults, FPGA with \code{rtl/fpga/*.sv}, Sky130 and ASAP7
  \code{__pnr__}.
\end{itemize}

\subsection{Differentiators (measured)}\label{tm:dec:xezim:6-differentiators-measured}\origsec{6}

{\footnotesize\setlength{\tabcolsep}{3pt}

{
\begin{longtable}[]{@{}
  >{\raggedright\arraybackslash}p{\dimexpr(\linewidth-2\tabcolsep)*265/1000\relax}
  >{\raggedright\arraybackslash}p{\dimexpr(\linewidth-2\tabcolsep)*715/1000\relax}@{}}
\toprule\noalign{}
\begin{minipage}[b]{\linewidth}\raggedright
Item
\end{minipage} & \begin{minipage}[b]{\linewidth}\raggedright
Result
\end{minipage} \\
\midrule\noalign{}
\endhead
\bottomrule\noalign{}
\endlastfoot
\textbf{Sky130 functional GLS} & Committed \code{rv32i_cpu_top} netlist:
\code{x4=0x0c}, BEQ taken. \textbf{Control:} July Synlig netlist
(\code{995c487}): \code{x4=0x30}, BEQ not taken. \textasciitilde14 s,
2.9 GB. 2 of 6 programs; no 31-register matrix \\
\textbf{SDF} & \textbf{No valid result.} No SDF exists for the committed
netlist. On a different on-disk netlist/\allowbreak{}SDF pair, xezim
annotated 0 signals and left 30,420 timing-check entries unmatched;
cause unknown \\
\textbf{X-propagation} & 10/10 probes match 4-state LRM expectation (no
commercial reference available). Verilator is 2-state; with
\code{rand+reset+2}, seed 7 takes the opposite \code{if} branch.
\code{XEZIM_INIT_ZERO=1} zeroes uninitialised memory but leaves scalar
\code{logic} at X \\
\textbf{UPF} & Small DUT: clamp and corruption occur under \code{--upf},
neither without. \code{phase5_soc.upf} parses and binds on
\code{soc_top}, but all 4 domains show "always on", so \textbf{no clamp
exercised}. Wrong supply path logs "unknown supply" and proceeds \\
\textbf{Metastability / CDC} & Nothing documented. Toy
\texttt{\$setuphold} + notifier gives X on violation;
\code{+notimingcheck} removes it \\
\end{longtable}
}

}

\subsection{Pros and cons per option}\label{tm:dec:xezim:7-pros-and-cons-per-option}\origsec{7}

{\footnotesize\setlength{\tabcolsep}{3pt}

{
\begin{longtable}[]{@{}
  >{\raggedright\arraybackslash}p{\dimexpr(\linewidth-6\tabcolsep)*93/1000\relax}
  >{\raggedright\arraybackslash}p{\dimexpr(\linewidth-6\tabcolsep)*360/1000\relax}
  >{\raggedright\arraybackslash}p{\dimexpr(\linewidth-6\tabcolsep)*360/1000\relax}
  >{\raggedright\arraybackslash}p{\dimexpr(\linewidth-6\tabcolsep)*168/1000\relax}@{}}
\toprule\noalign{}
\begin{minipage}[b]{\linewidth}\raggedright
Option
\end{minipage} & \begin{minipage}[b]{\linewidth}\raggedright
Pros
\end{minipage} & \begin{minipage}[b]{\linewidth}\raggedright
Cons
\end{minipage} & \begin{minipage}[b]{\linewidth}\raggedright
Best role
\end{minipage} \\
\midrule\noalign{}
\endhead
\bottomrule\noalign{}
\endlastfoot
\textbf{Verilator} & Fast build-once runs, in CI, mature, 896-test floor
& 2-state hides X, accepts use-before-declare, no SDF, zero-delay CDC &
Primary, unchanged \\
\textbf{VCS} & 4-state, full SV, 872/0/46 (reported), coverage works &
Licensed, not in CI, not on this host & Periodic cross-check \\
\textbf{Xcelium} & Same; fastest reported (\textasciitilde6 min SoC);
found \code{dft_clk_switch} race & Licensed; no \code{imc} for coverage
& Periodic cross-check \\
\textbf{xezim} & Free, 4-state event-driven, CI-capable, sub-second
elaboration, no build, GLS + UPF + X on one command & Silent miscompile
above; cocotb 1.9.2 blocked; 8 months old, one dominant author;
conflicting perf and sv-tests claims (below); no merge tool, FSM or
condition coverage; no Liberty & Local cross-check only \\
\textbf{Icarus} & Free, 4-state & Older SV; repo memory forbids it;
unused by any workflow & Remove \\
\textbf{Yosys \code{sim}} & Race-free cycle model, proved useful in
\code{u99} & Niche; netlist only & Keep for GLS triage \\
\textbf{sv2v + Yosys} & Sky130 fabric flow & Silent on missing modules;
slow on SoC (yosys \textasciitilde10 min) & Unchanged, synthesis only \\
\textbf{Synlig} & CPU macro baseline & Proven silent struct-select
defects & Guarded, per existing gates \\
\textbf{pyslang} & Strictest measured, fastest, already pinned &
Elaboration only; \code{--ignore-unknown-modules} trap &
\textbf{Blocking CI gate} (\code{lvdk}) \\
\textbf{SpyGlass / HAL} & Deep lint & Local only; CDC/RDC core rules not
run & Unchanged \\
\end{longtable}
}

}

Xezim claims not supported by this work: README near-parity with
commercial sim vs \code{NOTES.md:3288} "4-5x faster" reference; README
sv-tests 99.0 \% (0.11.0) vs checked-in 91.3 \% (0.8.1); README "keep
warm cache on" vs \code{--help} "EXPERIMENTAL, off by default";
\code{Cargo.toml} "MIT OR Apache-2.0" vs \code{LICENSE} Apache-2.0.

\subsection{Improvement areas}\label{tm:dec:xezim:8-improvement-areas}\origsec{8}

\textbf{This repo (independent of xezim):}

\begin{enumerate}
\def\labelenumi{\arabic{enumi}.}
\tightlist
\item
  Pyslang strict-elaboration gate (section 5): done, bead \bead{lvdk}.
  Closes the use-before-declare CI gap.
\item
  \src{uv.lock} pins cocotb 2.0.1 while \src{requirements.txt} says
  \code{<2.0}.
\item
  \texttt{sim/\allowbreak{}Makefile:\allowbreak{}90-\allowbreak{}91,\allowbreak{}\ 124-\allowbreak{}125}
  list \code{mem_collar.sv} and \code{mbist_ctrl.sv} twice; yosys fails
  on sv2v output unless de-duplicated (those two lists may serve
  different targets: check before editing).
\item
  Stale docs:
  \texttt{Doc:\allowbreak{}\ um:\allowbreak{}flow:\allowbreak{}ci-\allowbreak{}cd-\allowbreak{}workflow}
  (line 533 of the original file) suggests Icarus;
  \texttt{Doc:\allowbreak{}\ ds:\allowbreak{}ch:\allowbreak{}verification-\allowbreak{}plan\ (line\ 402\ of\ the\ original\ file)}
  says "Verilator + SDF", which is not implemented. Icarus remains in
  \code{flake.nix:203} + \src{tb/cocotb/cpu/Makefile:101}.
\item
  No free 4-state simulation in CI; X and race bugs surface only on
  non-CI proprietary runs. No clean solution today: xezim is the
  candidate, once fixed.
\item
  No local Verilator wall-clock baseline. CI timeouts
  aren\textquotesingle t measured durations.
\item
  Simulator choice hard-coded in \code{sim_portability.mk:56} and
  \code{run_cocotb_targets.py:152}; adding any simulator touches both.
\item
  Plan a cocotb 2 migration: \code{units=} was the recorded blocker, and
  it would also remove the 1.9.2 blocker here.
\end{enumerate}

\textbf{Xezim, if pursued (upstream, needs approval to file):} fix
section 4; settle cocotb 1.x re-entrancy; add cocotb to its CI (20
pytest cases currently not in \texttt{cargo\ test}); publish current
sv-tests data + a controlled benchmark; add
duplicate/\allowbreak{}excess override checks and opt-in read check;
report SDF mismatches clearly.

\subsection{Recommendation and open decisions}\label{tm:dec:xezim:9-recommendation-and-open-decisions}\origsec{9}

{\small\setlength{\tabcolsep}{4pt}

{
\begin{longtable}[]{@{}
  >{\raggedright\arraybackslash}p{\dimexpr(\linewidth-2\tabcolsep)*312/1000\relax}
  >{\raggedright\arraybackslash}p{\dimexpr(\linewidth-2\tabcolsep)*668/1000\relax}@{}}
\toprule\noalign{}
\begin{minipage}[b]{\linewidth}\raggedright
Role
\end{minipage} & \begin{minipage}[b]{\linewidth}\raggedright
Recommendation
\end{minipage} \\
\midrule\noalign{}
\endhead
\bottomrule\noalign{}
\endlastfoot
cocotb regression / CI & \textbf{Not now.} Re-check after section 4 is
fixed and either cocotb blocker clears \\
Local GLS and X cross-check & \textbf{Allowed as cross-check only}, with
a control able to fail. Disagreements are triaged, not trusted \\
Compile / elaboration check & \textbf{Use pyslang instead} \\
\end{longtable}
}

}

Open human decisions:

\begin{enumerate}
\def\labelenumi{\arabic{enumi}.}
\tightlist
\item
  \textbf{Decided:} pyslang gate added, blocking, on 6 arms
  (\src{tools/verif/check_elab_strict.py}, bead \bead{lvdk}).
\item
  \textbf{Done:} both reports filed upstream on 2026-11-02, after both
  reproduced on git \code{0f4d2fcf}: the section 4 miscompile as
  \href{https://github.com/aionhw/xezim/issues/810}{aionhw/\allowbreak{}xezim\#810}
  and the cocotb 1.9.2 abort as
  \href{https://github.com/aionhw/xezim/issues/811}{\#811} (bead
  \bead{gxg0}). Upstream fixed both the same day (\code{9cad2c17}).
  Measured on \code{d20039e4}: both testcases pass, and
  \code{test_mbist_soc_slice} passes 3/3 on stock cocotb 1.9.2 and on
  2.0.1, with no shim. All other results in this document are still from
  \code{9dcffd69}.
\item
  \textbf{Deferred:} xezim-core prototype, revisit after upstream
  responds (bead \bead{1ljh}).
\item
  \textbf{Open:} Phase 3 was deferred until the section 4 miscompile was
  fixed and a cocotb route existed. Both conditions are now met on
  \code{d20039e4}, so running it is a user decision (bead \bead{9xn3}).
\end{enumerate}

\subsection{What was NOT verified}\label{tm:dec:xezim:10-what-was-not-verified}\origsec{10}

\begin{itemize}
\tightlist
\item
  Full regression, SoC-scale cocotb suites, long-run throughput, memory
  on long runs.
\item
  Any cocotb 2 run of \code{test_gpio} or a pyuvm suite; size of cocotb
  2 migration.
\item
  VCS/\allowbreak{}Xcelium on these probes (use-before-declare confirmed
  by user, logs not seen here).
\item
  SDF on the committed macro; timing GLS of any kind; 4 of 6 GLS
  programs; register matrix; Verilator GLS timing.
\item
  Clamp and corruption on \code{soc_top}; X state on \code{soc_top}
  after reset beyond a sampled testbench (26 X-warn lines on MBIST,
  collar, macro signals; not triaged).
\item
  Root cause of section 4 and other expression forms; shim correctness
  beyond these suites; seeds other than 1234.
\item
  Any xezim source fix: section 2 fix assessments are from reading code.
\item
  Xezim\textquotesingle s own regression suite wasn\textquotesingle t
  run; its claimed C906/\allowbreak{}C910/\allowbreak{}Ibex results
  weren\textquotesingle t reproduced.
\end{itemize}
\end{relocated}

\begin{relocated}{docs/AGENT_IMPROVEMENT_PROPOSAL.md}
% Converted from docs/AGENT_IMPROVEMENT_PROPOSAL.md by tools/docs/md2tex.py; edit this file, not the original.
\section{\texorpdfstring{Proposed improvements to
\code{digital-chip-design-agents}}{Proposed improvements to }}\label{tm:dec:agent-improvement-proposal}

\textbf{Upstream:}
\href{https://github.com/hdl-tools/digital-chip-design-agents}{https:/\allowbreak{}/\allowbreak{}github.com/\allowbreak{}hdl-tools/\allowbreak{}digital-chip-design-agents}
\textbf{Versions audited:} all orchestrators at \code{1.2.0}
(\code{chip-design-memory-ip} is \code{1.8.0}, not audited)
\textbf{Basis:} 7 real agent invocations in one session on this repo
(2026-09-26/\allowbreak{}27), Phase 6a GPIO work. \textbf{Status:}
PROPOSAL. Nothing upstream has been changed.

\subsection{What happened}\label{tm:dec:agent-improvement-proposal:what-happened}

{
\begin{longtable}[]{@{}
  >{\raggedright\arraybackslash}p{\dimexpr(\linewidth-4\tabcolsep)*315/1000\relax}
  >{\raggedright\arraybackslash}p{\dimexpr(\linewidth-4\tabcolsep)*239/1000\relax}
  >{\raggedright\arraybackslash}p{\dimexpr(\linewidth-4\tabcolsep)*426/1000\relax}@{}}
\toprule\noalign{}
\begin{minipage}[b]{\linewidth}\raggedright
Agent
\end{minipage} & \begin{minipage}[b]{\linewidth}\raggedright
Task
\end{minipage} & \begin{minipage}[b]{\linewidth}\raggedright
Outcome
\end{minipage} \\
\midrule\noalign{}
\endhead
\bottomrule\noalign{}
\endlastfoot
\code{rtl-design-orchestrator} & GPIO RTL + SoC integration & Good;
needed 3 corrections from the caller \\
\code{verification-orchestrator} & SoC-level GPIO test & \textbf{Strong}
--- found a stale-build-cache hazard itself \\
\code{verification-orchestrator} & Debug an IRQ-clear failure &
\textbf{Strong} --- instrumented, root-caused, verified \\
\code{physical-design-orchestrator} & Hold-violation triage &
\textbf{Strong} --- corrected a bead\textquotesingle s false premise \\
\code{physical-design-orchestrator} & Slew/cap triage & \textbf{Strong}
--- \emph{declined to predict} beyond the evidence \\
\textbf{\code{firmware-orchestrator}} & \textbf{GPIO driver} &
\textbf{Failed} --- see below \\
\end{longtable}
}

\code{firmware-orchestrator} wrote a largely correct driver, then
\textbf{stopped without running its verification to completion}. Its
final report was \emph{"I\textquotesingle ll stop issuing commands now
and wait for the Monitor notification to arrive before continuing"} ---
after \textasciitilde250 k tokens spent looping on a wait for its own
background job. \textbf{The suite it produced failed when run, and the
agent never discovered this.} It also silently skipped the last item on
its deliverable list.

\subsection{The finding that matters most}\label{tm:dec:agent-improvement-proposal:the-finding-that-matters-most}

I expected the audit to show firmware lacking guards the strong
performers had. It does lack some (§1). But the bigger result is
\textbf{family-wide}:

\begin{quote}
\textbf{No orchestrator --- verification, PD, RTL or firmware ---
contains any instruction to verify before reporting, to quote exact
command output, or forbidding a claim that an unrun gate passed.}
\end{quote}

The strong performers succeeded \textbf{despite} their definitions, not
because of them. Three of four also omitted the \emph{same} final
deliverable (a CI pass-count floor bump) across unrelated tasks --- a
systemic completeness gap, not individual carelessness.

So the highest-value fixes are template-wide, not firmware-only.

\begin{center}\rule{0.5\linewidth}{0.5pt}\end{center}

\subsection{\texorpdfstring{Guards \code{firmware-orchestrator}
uniquely lacks}{Guards  uniquely lacks}}\label{tm:dec:agent-improvement-proposal:1-guards-firmware-orchestrator-uniquely-lacks}\origsec{1}

Measured presence across the four \code{1.2.0} orchestrators:

{\small\setlength{\tabcolsep}{4pt}

{
\begin{longtable}[]{@{}
  >{\raggedright\arraybackslash}p{\dimexpr(\linewidth-8\tabcolsep)*614/1000\relax}
  >{\centering\arraybackslash}p{\dimexpr(\linewidth-8\tabcolsep)*69/1000\relax}
  >{\centering\arraybackslash}p{\dimexpr(\linewidth-8\tabcolsep)*59/1000\relax}
  >{\centering\arraybackslash}p{\dimexpr(\linewidth-8\tabcolsep)*59/1000\relax}
  >{\centering\arraybackslash}p{\dimexpr(\linewidth-8\tabcolsep)*178/1000\relax}@{}}
\toprule\noalign{}
\begin{minipage}[b]{\linewidth}\raggedright
Guard
\end{minipage} & \begin{minipage}[b]{\linewidth}\centering
verif
\end{minipage} & \begin{minipage}[b]{\linewidth}\centering
PD
\end{minipage} & \begin{minipage}[b]{\linewidth}\centering
RTL
\end{minipage} & \begin{minipage}[b]{\linewidth}\centering
firmware
\end{minipage} \\
\midrule\noalign{}
\endhead
\bottomrule\noalign{}
\endlastfoot
Execution-tier hierarchy (MCP → wrapper → direct) & \ok{} & \ok{} &
\ok{} & \textbf{\bad{} no such section} \\
"These flows are long-running; read \code{metrics.json}/logs after
completion" & \bad{} & \ok{} & \bad{} & \bad{} \\
\texttt{Never\ proceed\ past\ a\ FAIL\ without\ applying\ the\ loop-\allowbreak{}back\ rule}
& \bad{} & \ok{} & \bad{} & \bad{} \\
\texttt{SUSPEND;\ flag\ RTL\ fix\ needed} terminal row & \ok{} & \bad{}
& \bad{} & \bad{} \\
\texttt{Escalate\ clearly\ if\ max\ iterations\ exceeded\ —\ show\ state\ and\ root\ cause}
& \bad{} & \bad{} & \ok{} & \bad{} \\
\end{longtable}
}

}

\code{firmware-orchestrator} is the \textbf{only} orchestrator with no
execution-tier section at all, and the only one with none of the five
guards above. It is also the shortest (82 lines vs
90/\allowbreak{}102/\allowbreak{}91).

\textbf{Especially incoherent:} it \emph{requires}
\texttt{stress\_\allowbreak{}test\_\allowbreak{}24h\_\allowbreak{}clean:\allowbreak{}\ true}
and a "24-hour high-throughput" stress test while giving \textbf{zero}
guidance on executing a long job. PD --- whose flows are obviously long
--- is the only one that got that guidance. Firmware demands a longer
job and got none.

\subsection{An outright defect: the experience log hardcodes success}\label{tm:dec:agent-improvement-proposal:2-an-outright-defect-the-experience-log-hardcodes-success}\origsec{2}

\code{chip-design-firmware/1.2.0/agents/firmware-orchestrator.md:78}
writes the session record with:

\begin{lstlisting}[language={}]
"signoff_achieved": true
\end{lstlisting}

hardcoded in the template, while the surrounding instruction says to
write the record \emph{"after signoff \textbf{or on
escalation/\allowbreak{}abandon}"}. So a failed or abandoned run appends
a record claiming success --- actively poisoning the
\code{memory/firmware/knowledge.md} corpus these agents are designed to
learn from. \textbf{This is a one-word fix and should land regardless of
everything else here.}

\begin{center}\rule{0.5\linewidth}{0.5pt}\end{center}

\subsection{Proposed changes}\label{tm:dec:agent-improvement-proposal:proposed-changes}

\subsubsection{\texorpdfstring{P1 --- Add a Reporting Contract to every
orchestrator
\emph{(template-wide)}}{P1 --- Add a Reporting Contract to every orchestrator (template-wide)}}\label{tm:dec:agent-improvement-proposal:p1-add-a-reporting-contract-to-every-orchestrator-template-wide}

Root cause of the only real failure: nothing obliged the agent to run
its gates or to prove it had.

\begin{lstlisting}[language={}]
## Reporting Contract
1. Before reporting, RUN every gate named in the task and paste each one's exact output.
2. Never report a gate as passing that you did not run in this session. If you could not run
   it, say so explicitly and say why.
3. A tool exiting 0 with empty or unparsable output is NOT a pass.
4. Re-read the task's deliverable list immediately before finishing and confirm each item.
   Report any you did not complete, and why.
5. Separate measured from inferred. Quote the number you observed; mark anything else as
   inference.
\end{lstlisting}

Item 3 is not hypothetical: this repo already hit it (bead \bead{dwp})
and wrote wrapper scripts to defend against it. Item 4 addresses the
3-of-4 missed-deliverable pattern. Item 5 is what the PD agent did
\emph{unprompted} --- it declined to predict a slew/cap regression the
evidence didn\textquotesingle t support --- and is worth making explicit
so it isn\textquotesingle t luck.

\subsubsection{\texorpdfstring{P2 --- Add a Long-Running Job protocol
\emph{(template-wide; generalise
PD\textquotesingle s)}}{P2 --- Add a Long-Running Job protocol (template-wide; generalise PD\textquotesingle s)}}\label{tm:dec:agent-improvement-proposal:p2-add-a-long-running-job-protocol-template-wide-generalise-pds}

Directly prevents the 250 k-token wait loop.

\begin{textblock}
## Long-Running Jobs
Builds, simulations and PD flows routinely exceed a single turn.
1. Start them in the background; never busy-poll in a loop.
2. Check at intervals matched to the job (minutes, not seconds). A quiet log during a
   compile is normal, not a hang — confirm liveness by CPU use before concluding anything.
3. If the job will outlive your turn budget, STOP and report: what is running, its job id,
   where its log is, and exactly what remains. A partial report naming the job is far more
   useful than an unverified success claim.
4. Never report results you have not seen. "Still running" is a valid, useful answer.
\end{textblock}

\subsubsection{\texorpdfstring{P3 --- Give \code{firmware-orchestrator}
the missing sections}{P3 --- Give  the missing sections}}\label{tm:dec:agent-improvement-proposal:p3-give-firmware-orchestrator-the-missing-sections}

Add the execution-tier hierarchy the other three have, plus
PD\textquotesingle s read-the-logs-after-completion rule, plus
RTL\textquotesingle s
\texttt{Escalate\ clearly\ …\ show\ state\ and\ root\ cause}, plus
PD\textquotesingle s \texttt{Never\ proceed\ past\ a\ FAIL}. Raise
\texttt{maxTurns:\allowbreak{}\ 70} toward PD\textquotesingle s
\code{100} \textbf{only after} P2 lands --- more turns without a wait
protocol just buys a longer stall.

\subsubsection{\texorpdfstring{P4 --- Fix the hardcoded
\code{signoff_achieved}}{P4 --- Fix the hardcoded }}\label{tm:dec:agent-improvement-proposal:p4-fix-the-hardcoded-signoff-achieved}

Replace the literal \code{true} at line 78 with the actual outcome, and
state that an abandoned or failed run must record \code{false}.

\subsubsection{\texorpdfstring{P5 --- Add an artifact-provenance guard
\emph{(new; from a defect that reached a
PR)}}{P5 --- Add an artifact-provenance guard (new; from a defect that reached a PR)}}\label{tm:dec:agent-improvement-proposal:p5-add-an-artifact-provenance-guard-new-from-a-defect-that-reached-a-pr}

No agent --- and neither did I --- asked whether \textbf{CI} could
obtain the generated artifacts a test depends on. A driver test consumed
\code{sw/bench/build/*.hex}, produced by a cross toolchain into a
\textbf{gitignored} directory. It passed locally and failed in CI with
\code{FileNotFoundError}.

\begin{textblock}
6. If a test consumes a generated artifact, verify its PROVENANCE in every environment that
   will run the test — not just yours. Check whether the artifact is committed, or
   reproducible by a step that environment actually performs. Passing locally because you
   built it by hand is not evidence CI will pass.
\end{textblock}

\begin{center}\rule{0.5\linewidth}{0.5pt}\end{center}

\subsection{Priority}\label{tm:dec:agent-improvement-proposal:priority}

{
\begin{longtable}[]{@{}
  >{\raggedright\arraybackslash}p{\dimexpr(\linewidth-8\tabcolsep)*54/1000\relax}
  >{\raggedright\arraybackslash}p{\dimexpr(\linewidth-8\tabcolsep)*338/1000\relax}
  >{\raggedright\arraybackslash}p{\dimexpr(\linewidth-8\tabcolsep)*160/1000\relax}
  >{\raggedright\arraybackslash}p{\dimexpr(\linewidth-8\tabcolsep)*80/1000\relax}
  >{\raggedright\arraybackslash}p{\dimexpr(\linewidth-8\tabcolsep)*348/1000\relax}@{}}
\toprule\noalign{}
\begin{minipage}[b]{\linewidth}\raggedright
\#
\end{minipage} & \begin{minipage}[b]{\linewidth}\raggedright
Change
\end{minipage} & \begin{minipage}[b]{\linewidth}\raggedright
Scope
\end{minipage} & \begin{minipage}[b]{\linewidth}\raggedright
Effort
\end{minipage} & \begin{minipage}[b]{\linewidth}\raggedright
Value
\end{minipage} \\
\midrule\noalign{}
\endhead
\bottomrule\noalign{}
\endlastfoot
P4 & Un-hardcode \code{signoff_achieved} & firmware, 1 line & trivial &
stops corrupting the memory corpus \\
P1 & Reporting Contract & all 4 & small & prevents shipping unverified
failures \\
P2 & Long-Running Job protocol & all 4 & small & prevents the 250
k-token stall \\
P3 & Backfill firmware\textquotesingle s missing sections & firmware &
small & brings it to parity \\
P5 & Artifact-provenance guard & all 4 & small & catches the
local-pass/\allowbreak{}CI-fail class \\
\end{longtable}
}

P4 and P1 are the two I would land first.

\subsection{Caveats on this evidence}\label{tm:dec:agent-improvement-proposal:caveats-on-this-evidence}

One session, one repo, seven invocations --- a small sample, and this
repo is unusually caveat-heavy, which may itself have shaped agent
behaviour. The \code{firmware-orchestrator} failure is a single data
point; the \emph{absence} of the guards in §1 and the
\code{signoff_achieved} defect are static facts about the files and do
not depend on sample size. The three-of-four missed-deliverable pattern
is behavioural and would want confirming on other repos before being
treated as settled.
\end{relocated}
